\section{Gaussian 与自适应求积}
\label{sec:08-Numerical-Analysis/04-Numerical-Differentiation-and-Integration/03}

被积函数不总是一行公式。它可能是一次流体仿真的输出，一次昂贵的前向模拟，一次要排队等三十毫秒的远程调用。你的外层循环要算几十万个一维积分，摊到每个积分上的预算是三次函数求值。

上一节的工具在这个预算下能做到什么？三个等距点给出 Simpson 公式，对三次以内的多项式精确。想再准一点，只能加点——可是加点正是你付不起的。

问题的提法里其实藏着一个没被动过的自由度。Newton--Cotes 那一族公式先把节点钉在等距位置上，然后解出权重；权重是 \(n+1\) 个待定数，节点却被当成了常量。若把节点也放开，三次求值对应的是六个自由参数而不是三个。

\subsection{节点位置也是可以调的变量}

放开之后能换来多少，可以先看结果。

\begin{center}
\begin{tikzpicture}[x=1cm,y=1cm,font=\small]
  \node[right,font=\footnotesize,black!70] at (-0.1,4.15) {同样的函数求值次数，节点摆在哪里};
  \draw[black!45] (0,3.5) -- (5.4,3.5);
  \draw[black!45] (0,3.35) -- (0,3.65);
  \draw[black!45] (5.4,3.35) -- (5.4,3.65);
  \foreach \X in {0,2.7,5.4} \fill[black!80] (\X,3.5) circle (2.2pt);
  \node[right,font=\footnotesize,black!80] at (5.6,3.5) {等距三点};
  \draw[black!45] (0,2.85) -- (5.4,2.85);
  \draw[black!45] (0,2.7) -- (0,3.0);
  \draw[black!45] (5.4,2.7) -- (5.4,3.0);
  \foreach \X in {0.609,2.7,4.791} \fill[blue!75] (\X,2.85) circle (2.2pt);
  \node[right,font=\footnotesize,blue!75] at (5.6,2.85) {Gauss 三点};
  \draw[black!45] (0,2.2) -- (5.4,2.2);
  \draw[black!45] (0,2.05) -- (0,2.35);
  \draw[black!45] (5.4,2.05) -- (5.4,2.35);
  \foreach \X in {0.253,1.246,2.7,4.154,5.147} \fill[orange!85] (\X,2.2) circle (2.2pt);
  \node[right,font=\footnotesize,orange!90] at (5.6,2.2) {Gauss 五点};
  \node[right,font=\footnotesize,black!70] at (-0.1,1.6) {能被精确积掉的最高多项式次数};
  \draw[fill=black!30,draw=black!45] (0,1.02) rectangle (1.5,1.28);
  \node[right,font=\footnotesize,black!80] at (1.65,1.15) {\(3\)};
  \draw[fill=blue!30,draw=black!45] (0,0.57) rectangle (2.5,0.83);
  \node[right,font=\footnotesize,blue!75] at (2.65,0.7) {\(5\)};
  \draw[fill=orange!35,draw=black!45] (0,0.12) rectangle (4.5,0.38);
  \node[right,font=\footnotesize,orange!90] at (4.65,0.25) {\(9\)};
\end{tikzpicture}

\emph{上半部分是三套节点在同一个区间上的位置，下半部分是各自的代数精确度。三个点不动、只把位置从等距挪到 Gauss 点，精确度从三次升到五次；五个 Gauss 点到九次，而五个等距点（Boole 公式）仍旧只有五次。注意 Gauss 节点全部落在区间内部，且三点与五点两套除中心外没有一个重合。}
\end{center}

对一个正的权函数 \(w\)，Gauss 求积写成

\[
\int_a^bw(x)f(x)\,dx\approx\sum_{i=1}^n\omega_if(x_i),
\]

节点取为关于 \(w\) 的 \(n\) 次正交多项式的根，权重全为正。\(n\) 个点对次数不超过 \(2n-1\) 的多项式精确——这是任何 \(n\) 点公式所能达到的上限。图里那三根长短不一的条，量的就是这件事。

\subsection{正交多项式的根是那组最优落点}

上限先说清楚。任何一组给定的 \(n\) 个节点，取

\[
g(x)=\prod_{i=1}^n(x-x_i)^2,
\]

它是 \(2n\) 次多项式，在所有节点上取零，求积和必为零；而 \(g\ge0\) 且不恒为零，对正权函数它的真实积分为正。所以 \(2n\) 次是过不去的。

达到 \(2n-1\) 的构造只有一行。记 \(\pi_n\) 为关于 \(w\) 的 \(n\) 次正交多项式，任意 \(p\in P_{2n-1}\) 做带余除法：

\[
p=q\pi_n+r,\qquad q,r\in P_{n-1}.
\]

积分那边，\(q\pi_n\) 的积分因正交性为零（\(q\) 的次数低于 \(n\)）；求积和那边，节点取在 \(\pi_n\) 的根上，\(\pi_n(x_i)=0\)，这部分也为零。两边只剩 \(r\)，而 \(r\) 的次数不超过 \(n-1\)，让权重对 \(P_{n-1}\) 精确即可——那正是插值型求积公式的定义。

条件对账：正交性防的是积分那一侧留下 \(q\pi_n\) 的残项，要求 \(q\) 的次数严格低于 \(n\)，所以 \(p\) 的次数不能超过 \(2n-1\)；根落在开区间内部（正交多项式的根都是实的、单的、在支撑内）防的是节点跑到积分区间之外；权函数为正防的是权重出现负号——负权重会在求和中制造相消，把上一单元里那种有效位损失重新引进来。

区间 \([-1,1]\)、\(w=1\) 时正交多项式是 Legendre 多项式。\(P_2(x)=\tfrac12(3x^2-1)\) 的根给出两点公式 \(x_{1,2}=\pm1/\sqrt3\)，权重都是 \(1\)：

\[
\int_{-1}^1f(x)\,dx\approx f(-1/\sqrt3)+f(1/\sqrt3),
\]

两次求值精确积掉全部三次多项式。搬到一般区间 \([a,b]\) 要做仿射变换，节点按 \(x\mapsto\frac{a+b}{2}+\frac{b-a}{2}x\) 平移伸缩，权重同时乘 \(\frac{b-a}{2}\)；只换节点不换权重会让积分整体差一个尺度因子，是这类代码最常见的一处错。

权函数是可以挑的。把 \(e^{-x^2}\) 挑进 \(w\) 得到 Gauss--Hermite，把 \(e^{-x}\) 挑进去得到 Gauss--Laguerre，把端点奇性 \((1-x)^\alpha(1+x)^\beta\) 挑进去得到 Gauss--Jacobi。规则与测度对上号时，被积函数里最难缠的那部分被吸收进权重表，剩下的部分光滑，少数几个节点就够。计算高斯型期望时这一点特别值钱：\(\E[g(Z)]\) 对 \(Z\sim\mathcal N(0,1)\)，用五个 Hermite 节点常常比一万个随机样本还准。

Gauss 节点向区间两端加密的形态，与\hyperref[sec:08-Numerical-Analysis/03-Interpolation-and-Approximation/02]{``Runge 现象与节点选择''}里的 Chebyshev 节点是同一副长相。两处的理由也相通：等距布点在端点附近让节点乘积鼓出巨包，把节点往两端挪能把这个包压平。求积和插值在这件事上共用一个答案。

\subsection{代价是这组节点不能复用}

精度不是白拿的。Gauss 节点是无理数，位置随 \(n\) 整个改变——图里三点与五点两排，除了中心那个点，没有一个能对上。这意味着算完三点想加密到五点，前三次求值全部作废，要重头再算五次。等距的复合公式没有这个问题：从 \(n\) 段加密到 \(2n\) 段，原有的函数值一个不丢，只需补上新增的中点，这也正是上一节 Richardson 外推能成立的前提。

Gauss--Kronrod 就是为这笔账设计的补丁。它在 \(n\) 个 Gauss 节点之外再嵌入 \(n+1\) 个新节点，凑成一个 \(2n+1\) 点的规则；两套规则共享那 \(n\) 个已有的函数值，差值直接当误差指标用。代价是嵌入规则的代数精确度低于同点数的纯 Gauss 规则——拿一部分精度换一个免费的误差估计。几乎所有工业级的一维积分例程内部跑的都是它。

\subsection{自适应把预算搬到函数难缠的那一段}

前面讨论的都是固定网格：节点在哪由公式决定，与被积函数长什么样无关。可被积函数的困难往往集中在很小的一块。

\begin{center}
\begin{tikzpicture}[x=1.7cm,y=1.5cm,font=\small,>=stealth]
  \draw[black!30] (-0.05,0) -- (4.2,0);
  \draw[black!45] (0,0) -- (0,1.55);
  \draw[black!85,very thick] plot[smooth] coordinates {
    (0,0.250) (0.8,0.250) (1.5,0.250) (2.0,0.250) (2.1,0.255) (2.2,0.283)
    (2.3,0.406) (2.4,0.723) (2.5,1.171) (2.6,1.400) (2.7,1.171) (2.8,0.723)
    (2.9,0.406) (3.0,0.283) (3.1,0.255) (3.2,0.250) (3.6,0.250) (4.0,0.250)};
  \draw[black!35,dashed] (2.2,0) -- (2.2,1.45);
  \draw[black!35,dashed] (3.0,0) -- (3.0,1.45);
  \foreach \t in {0,0.25,...,4} \fill[black!70] (\t,-0.30) circle (1.4pt);
  \foreach \t in {0,0.6,1.2,1.8,2.1,2.2,2.3,2.375,2.45,2.525,2.6,2.675,2.75,2.825,2.9,3,3.2,3.6,4}
    \fill[blue!75] (\t,-0.62) circle (1.4pt);
  \node[right,font=\footnotesize,black!75] at (4.15,-0.30) {均匀 \(17\) 点};
  \node[right,font=\footnotesize,blue!75] at (4.15,-0.62) {自适应 \(19\) 点};
  \node[font=\footnotesize,black!70] at (1.05,1.15) {平坦段};
  \node[font=\footnotesize,black!70] at (3.75,1.15) {平坦段};
\end{tikzpicture}

\emph{几乎相同的求值预算，两种花法。均匀网格把十三个点撒在两侧的平坦段上，落进虚线之间那个窄峰的只有四个；自适应网格在平坦段几点带过，把十个点压进峰区。峰的面积占全积分的大半，误差也几乎全在那里。}
\end{center}

自适应求积的机制很简单：在每个区间上算一粗一细两个近似，用两者之差估计局部误差；达标就收下，不达标就把区间劈成两半，对每一半重复。图里那种点的分布不是谁摆出来的，是这个递归自己长出来的。

误差估计可以借上一节的渐近关系。设整个区间的一次 Simpson 结果是 \(S\)，左右半区间之和是 \(S_L+S_R\)。四阶意味着步长减半误差降到十六分之一，于是 \(I-S\approx16(I-S_L-S_R)\)，解出

\[
I-(S_L+S_R)\approx\frac{(S_L+S_R)-S}{15}.
\]

右边全是已经算出来的量。这个 \(1/15\) 就是 Richardson 外推倒过来用：那边拿幂次关系去消掉误差，这边拿同一个幂次关系去读出误差。Gauss--Kronrod 的误差指标是另一种读法，用两套阶数不同的规则代替两种步长。

实现时还有一串必须显式设定的东西：绝对容差与相对容差各一个（只设相对容差会在积分值接近零时永不收敛）、最大求值次数、递归深度上限、最小区间宽度，以及 \(\infty\) 与 NaN 出现时的处置。这些不是工程琐碎，是防住递归在奇点处无限细分下去的唯一手段。同一套自适应逻辑在\hyperref[sec:08-Numerical-Analysis/06-ODE/02]{``Runge--Kutta 与自适应步长''}里控制的是时间步长，控制器的骨架完全一样。

误差估计本身也会骗人。粗细两个结果接近，只说明加密这一次没带来变化，不说明结果对。若两套节点都从一个极窄的峰旁边跨过去，差值会很小，算法安心地收下一个错误答案——这与上一节比较 \(T_n\) 与 \(T_{2n}\) 时的陷阱是同一个。防它的办法不在数值方法内部：知道峰在哪就手动拆区间，或者先粗扫一遍定位可疑区域。

\subsection{不光滑时高阶反而吃亏}

代数精确度描述的是一个测试空间：规则在多项式上表现如何。被积函数离多项式很远时，这个指标失去参考价值。

Gauss 规则的误差里带着 \(f^{(2n)}\)。函数光滑、高阶导数有界时，这一项随 \(n\) 迅速衰减，收敛快得惊人；解析函数上甚至是几何速率，这条线索在\hyperref[sec:08-Numerical-Analysis/08-Fourier-and-Spectral-Methods/04]{``谱方法：用光滑性换高精度''}里还会展开。可一旦函数在区间内有跳跃、有尖点、有端点奇性，\(f^{(2n)}\) 根本不存在，高阶规则不再有任何优势，甚至比低阶更糟——它在更宽的邻域内假定一个高次多项式能代表函数，而这个假定错得更远。

这种时候正确的动作是降阶加密，不是升阶。分段低次公式加上足够多的段，收敛速度由函数实际的光滑度决定，不会因为方法的野心而恶化。更好的做法是先把不光滑处摘出来：在跳跃点把区间拆开，让每一段内部光滑；对已知的端点奇性做变量替换或奇性减除；把奇性吸进权函数，交给匹配的 Gauss--Jacobi 规则。

选规则的时候，这几条大致够用。手上只有一张离散数据表、无法在任意点求值，那么分段梯形是最诚实的选择——你没有凭空使用未采样曲率的权利。函数光滑且可以随意取样，Gauss 或 Clenshaw--Curtis 用很少的点就能到高精度。困难集中在局部，交给自适应。被积函数是解析的周期函数且积分覆盖整周期，复合梯形法反而异常优秀，因为误差展开中所有带端点导数的项两两抵销。端点奇性或无限区间，先做变换或换匹配的权函数，再谈阶数。

而一维的所有讨论到高维都要重来。张量积把 \(n\) 点公式变成 \(n^d\) 点，维数灾难直接落在求值次数上，稀疏网格能缓解一部分，超过十几维就得换到 Monte Carlo 那一类不依赖光滑度、收敛率与维数无关的方法上去。求积规则的精妙全建立在``函数在低维网格上可被多项式代表''这个前提上，前提不成立时，精妙本身就是负担。
